\begin{table}[H]
\centering
\caption{Summary Statistics: NAICS 512 Employment by Treatment Status}
\label{tab:summary}
\begin{threeparttable}
\begin{tabular}{lcccc}
\toprule
 & \multicolumn{2}{c}{Treated States} & \multicolumn{2}{c}{Never-Treated} \\
\cmidrule(lr){2-3} \cmidrule(lr){4-5}
 & Mean & SD & Mean & SD \\
\midrule
Employment (all races) & 7592 & 13574 & 22376 & 75270 \\
Hires (all races) & 3403 & 7178 & 8517 & 29038 \\
Earnings (\$) & 1498 & 927 & 1268 & 1052 \\
\addlinespace
\multicolumn{5}{l}{\textit{Treated states, by race:}} \\
\quad White employment & 6030 & 10495 & & \\
\quad Black employment & 999 & 2032 & & \\
\quad Hispanic employment & 21 & 61 & & \\
\addlinespace
States & 37 & & 14 & \\
State-quarter obs & 3207 & & 1189 & \\
\bottomrule
\end{tabular}
\begin{tablenotes}[flushleft]\small
\item \textit{Notes:} Data from Census QWI (race/ethnicity table, NAICS 512 Motion Picture), quarterly 2003--2024. Treated states adopted film production tax credits ($\geq$15\% rate) between 2002 and 2019. Never-treated states include AK, CA, DE, ID, IA, KS, ND, NE, NH, SD, VT, WY, DC. Employment, hires, and earnings are state-quarter means.
\end{tablenotes}
\end{threeparttable}
\end{table}

